\begin{frame}[plain]
\wrsTitlePage{Diagram backends}{Which backend should draw which scientific visual?}{Week 1 \textperiodcentered{} one semantic figure spec, rendered by TikZ, Draw.io and native Beamer layout.}{Weekly research update \textperiodcentered{} synthetic demo \textperiodcentered{} version 0.1}
\note{Name the routing question. Do not explain each backend yet; the deck shows them.}
\end{frame}

\begin{frame}{What stays constant across renderers}
\begin{columns}[T]
\column{0.5\textwidth}
\wrsSection{Established}
\begin{itemize}
  \item \textbf{Method model} One semantic model of the method
  \item \textbf{Figure spec} Nodes
  \item \textbf{Routing} The renderer is chosen explicitly or by complexity heuristics.
\end{itemize}
\column{0.46\textwidth}
\wrsSection{This week}
\wrsBody{Four renderers coexist in one deck. Conceptual diagrams are native TikZ, the large architecture stays in Draw.io, results stay a native table, and the deck is still one Beamer PDF.}\par\vspace{0.4em}
\wrsTiny{Previous: separate figure pipelines}
\end{columns}
\note{Keep it short. The point is that the science does not move between renderers.}
\end{frame}

\begin{frame}{Direct reuse vs learned correction}
\wrsLead{The difference should be visible in under five seconds.}\par\vspace{0.5em}
\wrsDiagram[0.97\linewidth]{figures/cached_vs_corrected.tex}
\vspace{0.3em}\wrsQuiet{Same inputs, same cache, same output; only the operation between cache and output differs. TikZ keeps both panels in one coordinate system.}
\note{Point at the middle operation in both panels. Everything else is deliberately identical.}
\end{frame}

\begin{frame}{Correction direction vs desired direction}
\wrsDiagram{figures/preimage_geometry.tex}
\vspace{0.3em}\wrsQuiet{The correction is nearly orthogonal to the desired direction; the angle is annotated from the same measurement as the diagnostic slide.}
\note{Reveal objects first, then the desired direction, then the correction. The angle is the observation.}
\end{frame}

\begin{frame}{Method delta: v3 to v4}
\wrsDiagram[0.97\linewidth]{figures/v3_v4.tex}
\vspace{0.3em}\wrsQuiet{Unchanged components keep their layer across both panels and stay muted; the inserted correction is the only highlighted object.}
\note{Compare left and right by position, not by reading every label.}
\end{frame}

\begin{frame}{Large architecture stays editable in Draw.io}
\wrsFigure[0.92\linewidth]{assets/lesa_overview.pdf}
\vspace{0.3em}\wrsQuiet{Draw.io remains the backend for large editable architecture figures; Beamer includes the exported PDF directly.}
\note{This figure is intentionally not TikZ; the router would choose Draw.io for it.}
\end{frame}

\begin{frame}{Renderer cost and editability}
\begin{center}\scriptsize
\setlength{\tabcolsep}{3pt}
\renewcommand{\arraystretch}{1.15}
\begin{tabular}{lccc}
\toprule
Method & \textbf{Compile (s)} & \textbf{Native typography} & \textbf{Editable source} \\
\midrule
\textbf{\textcolor{wrsBlue}{TikZ conceptual}}\newline{\tiny\color{wrsMuted}(this week)} & \textbf{\textcolor{wrsBlue}{1.20\ {\tiny\color{wrsGreen}(once)}}} & \textbf{\textcolor{wrsBlue}{1\ {\tiny\color{wrsGreen}(yes)}}} & \textbf{\textcolor{wrsBlue}{0.50\ {\tiny\color{wrsGreen}(tex)}}} \\
Draw.io architecture\newline{\tiny\color{wrsMuted}(last week)} & 0.40\ {\tiny\color{wrsGreen}(pdf)} & 0.50\ {\tiny\color{wrsGreen}(via pdf)} & 1\ {\tiny\color{wrsGreen}(gui)} \\
Python plot & 0.30\ {\tiny\color{wrsGreen}(pre)} & 0.50\ {\tiny\color{wrsGreen}(no)} & 0.50\ {\tiny\color{wrsGreen}(script)} \\
\bottomrule
\end{tabular}
\end{center}
\vspace{0.25em}\wrsTiny{lower is better; 1 = native; 1 = full}
\vspace{0.25em}\wrsQuiet{TikZ costs a compile step and pays back in native typography and direct Beamer integration; Draw.io stays best for large editable figures.}
\note{Honest trade-off, not a benchmark of quality.}
\end{frame}

\begin{frame}{Limitations and open questions}
\begin{columns}[T]
\column{0.5\textwidth}
\wrsSection{Current limitations}
\begin{itemize}
  \item TikZ routing is heuristic; the decision is recorded but not yet learned from feedback.
  \item Only conceptual and diagnostic archetypes are implemented in v0.1.
  \item PGFPlots is optional and not used by the default plot pipeline.
  \item Overlays inside a scaled figure are not supported.
\end{itemize}
\column{0.46\textwidth}
\wrsSection{Open questions}
\begin{itemize}
  \item Should routing learn from critic outcomes across weeks?
  \item Can the same figure spec drive a paper-ready export automatically?
  \item When does a TikZ diagram become too large to stay readable in a slide?
\end{itemize}
\end{columns}
\note{Stop here. Discussion follows.}
\end{frame}
